% Transcription — fonds Alexandre Grothendieck, shelfmark 10, batch 1 (pages 1-20).
% Established from the facsimile archives/batches/10/batch-01.pdf (a mirror of
% https://grothendieck.umontpellier.fr/10.pdf), per the skill
% transcribe-grothendieck.
%
% Pass: Fable 5.1 (claude-fable-5-1), 2026-09-10 — first pass, unchecked against
% the pages by a human.
%
% The batch holds two runs. Pages 1-16 are one continuous draft in blue ink,
% opened by a cover sheet inscribed « Polarisations sur les ⊗-catégories
% graduées » (page 1, no \page{}), and numbered by him 1) to 19):
%
%   3-5   Weil forms on a ⊗-category over K ⊂ ℝ: the anti-involution u ↦ u*
%         attached to a non-degenerate ε-symmetric φ: M⊗M → L, the Weil
%         condition Tr u*u > 0, compatibility of two Weil forms, its failure
%         to be transitive, the computation u^{ad(φ,ψ)} = α^{-1}u*α, and the
%         definition of a prépolarisation P(M) ⊂ Hom(M⊗M, K(−i)) with (P1),
%         the maximal prépolarisation P̂ and its corollaries.
%   6-7   admissibility (P2), (P3), polarisation = maximal and admissible,
%         the six stability properties a)-f), the reduction to generators
%         ℳ_α and to ⊗-categories over ℝ (Remarques 1-3). Page 7 is a fast
%         draft whose connective prose is largely lost.
%   8-9   complexified objects M ∈ Ob ℳ^i_ℂ with J_M, hermitian and
%         anti-hermitian forms, φ_0 invariant under J_M, M̌ ≃ M̄(−i), and the
%         chain α) ⇒ β) ⇒ γ) with counter-examples in Rep(G).
%   10-11 the criterion for a polarisation of a semi-simple ℳ generated by
%         M_1…M_r with forms φ_i, through M_m and φ_m, and its sufficiency.
%   11-16 ℳ = Rep(G), G_m → G → G_m, η = i(−1), C ∈ G(ℝ) with C² = η, the
%         C-polarisations, the questions of 16), the centraliser of C as a
%         maximal compact of G'(ℝ) (17), when two C give the same
%         polarisation and the torsor under ₂Z(ℝ) (18), and the action of
%         ₂Aut_⊗(id_ℳ) on Pol(ℳ) in general (19).
%
% Pages 18 and 20 are a second, older run in black ink on yellowed paper,
% headed « Notes anciennes » in his hand (top right of page 18): the adjoint
% u ↦ u^φ of an endomorphism of a motive M of weight p relative to a form φ
% with values in Q(p), through Θ_φ: End(M) ≃ 𝒬(M) = B(M,M;Q(p)) ≃ Hom(M,M*),
% and the sign (−1)^δ.
%
% Page 2 is the back of the cover, page 17 a blank leaf, page 19 a blank
% leaf with one struck line upside down at its foot — all three absent.
%
% The dating is the inventory's own, for the whole shelfmark.
\documentclass[11pt,a4paper]{article}
% Le préambule commun à toutes les transcriptions du fonds Grothendieck.
%
% Il tient en une idée : **l'incertitude se compose, elle ne se résout pas.**
% Une transcription qui lisse un mot douteux a détruit la seule chose pour
% laquelle on la faisait — un lecteur doit pouvoir savoir, à chaque endroit,
% ce qui était sur le feuillet et ce qui a été ajouté. Les six macros qui
% suivent sont donc l'appareil critique, et non de la décoration.
%
% Les mêmes commandes sont reconnues par `scripts/render.mjs`, qui produit la
% vue de lecture du site. Une macro ajoutée ici doit l'être aussi là-bas,
% faute de quoi le rendu échouera bruyamment — ce qui est voulu : un rendu
% silencieusement faux est pire qu'un rendu absent.

\ProvidesPackage{grothendieck}[2026/08/08 Transcriptions du fonds Grothendieck]

\usepackage{amsmath,amssymb,amsthm}
\usepackage{mathrsfs}
\usepackage{stmaryrd}
% Les diagrammes commutatifs sont partout dans ce fonds : c'est la langue
% même de la géométrie algébrique de Grothendieck, et une transcription qui
% les rendrait en prose ne transcrirait rien.
\usepackage{tikz-cd}
% Français par défaut — c'est la langue des feuillets — mais l'anglais est
% chargé aussi : la traduction et le résumé sont des documents anglais, et la
% typographie française (espaces avant les deux-points) y serait une faute.
% Ces documents basculent par \selectlanguage{english} en tête de corps.
\usepackage[english,french]{babel}
\usepackage{fontspec}
\usepackage{geometry}
\geometry{a4paper,margin=25mm,marginparwidth=28mm}
\usepackage{marginnote}
\usepackage{xcolor}
% ulem, not soul. `soul`'s \st and \ul are built on a character-by-character
% scan that XeLaTeX's font machinery breaks: the document fails with an
% undefined control sequence rather than degrading. ulem does the same two jobs
% and is Unicode-engine safe. `normalem` stops it hijacking \emph, which the
% transcriptions use for Grothendieck's own emphasis.
\usepackage[normalem]{ulem}
\usepackage{hyperref}
\hypersetup{colorlinks=true,linkcolor=black,urlcolor=black,pdfborder={0 0 0}}

% --- Métadonnées du lot -----------------------------------------------------
% Reprises telles quelles par le script de rendu, qui les affiche en tête de
% la vue de lecture. Elles fixent ce que le fichier prétend couvrir : sans
% elles, une transcription flotte sans point d'attache dans un fonds de seize
% mille feuillets.

\newcommand{\folder}[1]{\gdef\grothFolder{#1}}
\newcommand{\batch}[1]{\gdef\grothBatch{#1}}
\newcommand{\leaves}[2]{\gdef\grothFirst{#1}\gdef\grothLast{#2}}
\newcommand{\foldertitle}[1]{\gdef\grothFoldertitle{#1}}
\gdef\grothFolder{?}\gdef\grothBatch{?}\gdef\grothFirst{?}\gdef\grothLast{?}\gdef\grothFoldertitle{}

% --- Le résumé --------------------------------------------------------------
%
% Il ouvre la lecture modernisée, sous le titre « Résumé », et il est composé
% en petit. Ce n'est pas un ornement : c'est le seul passage du document qui
% ne soit pas des mathématiques, et un lecteur doit pouvoir le reconnaître
% avant de l'avoir lu — puis le sauter s'il connaît déjà le sujet, ou s'y
% arrêter s'il ne le connaît pas. Le filet qui le suit marque la reprise du
% corps.

\newenvironment{resume}
  {\small\setlength{\parskip}{0.45em}}
  {\par\medskip\hrule\medskip}

% --- Mots-clés ---------------------------------------------------------------
%
% En anglais, en clôture du résumé de la lecture modernisée : le vocabulaire
% moderne sous lequel chercher ce que les feuillets construisent. Le manifeste
% du site (`npm run manifest`) les extrait pour étiqueter la cote entière —
% cette ligne est la source unique des tags, il n'y a pas de fichier à côté.
\newcommand{\keywords}[1]{\par\smallskip\noindent
  {\footnotesize\textsc{Keywords} --- \textit{#1}}\par}

% --- L'appareil critique ----------------------------------------------------

% Le numéro de feuillet, en marge. C'est l'ancre qui permet de régler un
% désaccord sur une formule en regardant *un* feuillet plutôt qu'une chemise
% de six cents. Le site s'en sert aussi pour tourner les pages du fac-similé
% quand on fait défiler la transcription.
\newcommand{\leaf}[1]{%
  \marginnote{\footnotesize\color{gray!70}\textsf{#1}}%
  \label{leaf:#1}\ignorespaces}

% Illisible : on ne devine pas. Un mot inventé qui se lit comme les autres est
% le pire résultat possible.
\newcommand{\ill}{\textcolor{red!60!black}{[\,\ldots\,]}}

% Lecture douteuse : le mot est donné, et signalé comme incertain.
\newcommand{\uncertain}[1]{\uline{#1}}

% Ajout de l'éditeur — une lettre manquante, un mot sous-entendu.
\newcommand{\add}[1]{\textcolor{blue!50!black}{[#1]}}

% Biffé par Grothendieck lui-même. Ce qu'il a rayé fait partie du document :
% une rature dit souvent où la pensée a changé de direction.
\newcommand{\struck}[1]{\sout{#1}}

% Note du transcripteur, sur l'état matériel du feuillet.
\newcommand{\note}[1]{\par\smallskip{\small\color{gray!60!black}\textit{[#1]}}\par\smallskip}

% Note marginale de Grothendieck, distincte du corps du texte.
\newcommand{\marginal}[1]{\par\smallskip\noindent\hspace{1em}%
  {\small\color{gray!60!black}#1}\par\smallskip}

% --- Titre ------------------------------------------------------------------

\AtBeginDocument{%
  \begin{center}
    {\large\bfseries Fonds Alexandre Grothendieck --- cote n\textsuperscript{o}~\grothFolder}\\[2pt]
    {\itshape\grothFoldertitle}\\[4pt]
    {\small feuillets \grothFirst--\grothLast\ (lot \grothBatch)}\\[2pt]
    {\footnotesize\color{gray!60!black}Transcription établie d'après le fac-similé numérisé
     par l'Université de Montpellier.\\
     Les lectures incertaines sont soulignées, les illisibles notées [\,\ldots\,],
     les ajouts entre crochets.}
  \end{center}
  \vspace{1em}}

\endinput


\folder{10}
\batch{1}
\pages{1}{20}
\dating{[à partir de 1958]}
\watermark{Édition de démonstration}
\foldertitle{Catégories tensorielles : notes manuscrites (s.d.), tapuscrits (s.d.)}

\begin{document}

\section*{Polarisations sur les $\otimes$-catégories graduées (pages 3 à 16)}

\note{Le titre est de la main de Grothendieck, sur la chemise qui précède
(page 1) : « Polarisations sur les $\otimes$-catégories graduées ». Les
paragraphes sont numérotés par lui de 1) à 19). Dans tout le feuillet,
« ½-simple » abrège \emph{semi-simple}, « ss-objet » \emph{sous-objet},
« comp. : » \emph{compatible avec}, « rep. » \emph{représentation}, et
$u^{*(\varphi,\psi)}$, $u^{\mathrm{ad}(\varphi,\psi)}$ désignent tous deux
l'adjoint de $u \colon M \to N$ relativement aux formes $\varphi$ sur $M$ et
$\psi$ sur $N$. L'indice $2$ à gauche, ${}_2Z(\mathbb{R})$, désigne les
éléments d'ordre divisant $2$.}

\page{3}
\struck{$\mathcal{M}$ catégorie abélienne $K$-linéaire, $K \subset \mathbb{R}$}
\note{Sous « $K \subset \mathbb{R}$ » biffé, deux mots minuscules, \ill{}.}

$\mathcal{M}$ une $\otimes$-catégorie sur $K \subset \mathbb{R}$, avec
\struck{\ill{}} objet inversible \struck{$K(1)$ de degré 2.} $L$

1) Soit $M \in \mathrm{Ob}\,\mathcal{M}$, $\varphi \colon M \otimes M \to L$.
\note{$\mathcal{M}$ porte un exposant serré, peut-être $i$, non lu.}
Supposons $\varphi$ symétrique ou alternée, $\varphi$ non dégénérée. Alors
$\varphi$ définit sur $\mathrm{End}(M)$ une \emph{anti-involution}
$u \mapsto u^{*}$ [$\varphi(ux, y) = \varphi(x, u^{*}y)$ ; posant
$\varphi(x,y) = \langle x, \tilde{\varphi} y\rangle$,
$\tilde{\varphi} \colon M \to \check{M} \otimes L$, avec donc
$u^{*} = \tilde{\varphi}^{-1} \circ \check{u} \circ \tilde{\varphi}$].
\struck{\ill{}} On a une bijection
\[
\mathrm{Hom}(M \otimes M, L) \simeq \mathrm{End}(M)
\]
[$\psi$ et $u$ se correspondent si $\psi(x,y) = \varphi(x, uy)$, i.e.\
$\tilde{\psi} = \tilde{\varphi} u$], et cette dernière \struck{bijection}
transforme $\psi \mapsto {}^{t}\psi$ en $u \mapsto \varepsilon u^{*}$, (où
$\varepsilon = \pm 1$ suivant que $\varphi$ sym.\ ou alternée). Donc les
formes sym.\ resp.\ alternées) de $\mathrm{Hom}(M \otimes M, L)$
correspondent aux $u$ tels que $u^{*} = \varepsilon u$.
\note{Le facteur $\varepsilon$ de la dernière égalité est d'une lecture douteuse.}

2) On dit que $\varphi$ (sym.\ ou altern.) \struck{est définie} \add{satisfait
la condition de Weil} si $\operatorname{Tr} u^{*}u > 0$ pour
$u \in \mathrm{End}(M)$, $u \neq 0$. (Ceci implique, si $\mathrm{End}(M)$ de
dim.\ finie sur $K$, que $\mathrm{End}(M)$ est ½-simple, et pour
$\forall$ ss-objet $N \hookrightarrow M$, $\varphi|N$ est une forme de Weil.)
\marginal{donc que $M$ est semi-simple, s'il est de \uncertain{long.\ finie}
et \uncertain{si tt ss-objet a un complément direct}}
\note{« satisfait la condition de Weil » est une insertion interlinéaire
au-dessus des mots biffés ; la note marginale est cerclée et reliée d'un
trait à la parenthèse.}

3) Soient $\varphi \in \mathrm{Hom}(M \otimes M, L)$,
$\psi \in \mathrm{Hom}(N \otimes N, L)$, $\varphi$ et $\psi$
$\varepsilon$-sym.\ ($\varepsilon = \pm 1$). Alors $\varphi \oplus \psi$
satisfait la condition de Weil ssi $\varphi$ et $\psi$ y satisfont, et si de
plus, pour $\forall\, u \colon M \to N$, on a $\operatorname{Tr} u^{*}u > 0$
si $u \neq 0$. [$u^{*}$ calculé pour les formes $\varphi$ et $\psi$ sur $M$ et
$N$]
\marginal{Généralisation au cas où on a \struck{\ill{}} un nombre quelconque
de \uncertain{facteurs} \uncertain{(ou sommandes)}}

\page{4}
On dit alors que les formes $\varphi$ et $\psi$ \struck{(satisfaisant aux
conditions de Weil)} sont \emph{compatibles}.

4) La compatibilité des formes de Weil est une relation \struck{d'équivalence}
réflexive et symétrique. \struck{\ill{}}
\note{Trois lignes biffées, dont on lit « Si $\varphi, \varphi'$ sur $M$,
$\psi, \psi'$ sur $N$ » et « compatibles » ; une note marginale, écrite en
biais dans la marge gauche et rattachée au passage biffé, dit pourquoi la
relation n'est pas transitive.}
\marginal{\struck{\ill{}} relation \ill{} transitive : \uncertain{si $M$, $N$
sont tels que $\mathrm{Hom}(M,N) = \mathrm{Hom}(N,M) = 0$}, $M \neq 0$,
$\varphi$, $\varphi'$ \struck{formes de Weil} sur $M$, $\psi$ sur $N$,
$\varphi$ compatible avec $\psi$ et $\psi$ compatible avec $\varphi'$,
\uncertain{mais pas $\varphi$ compatible avec $\varphi'$}}
Si $\varphi, \varphi'$ compatibles, $\psi, \psi'$ compatibles, alors
$\varphi$ et $\psi$ sont compatibles ssi $\varphi'$ et $\psi'$ le sont. Pour
le voir, il faut expliciter.

5) Soient $\varphi$ et $\psi$ deux formes $\varepsilon$-symétriques
$M \otimes M \to L$, $u \mapsto u^{*}$ l'anti-involution de $\mathrm{End}(M)$
associée à $\varphi$, et $\psi$ définie par $\alpha \in \mathrm{End}(M)$,
$\psi(x,y) = \varphi(x, \alpha y)$, i.e.\
$\tilde{\psi} = \tilde{\varphi} \circ \alpha$. Alors $u^{\mathrm{ad}(\psi)}$
défini par $\psi$ \struck{\ill{}} c'est : $\alpha^{-1} u^{*} \alpha$,
\struck{donc $\psi$ est compatible avec $\varphi$ ssi $\operatorname{Tr}$} et
$u^{\mathrm{ad}(\varphi,\psi)}$ \struck{défini par $\psi$} c'est :
$u^{*}\alpha$. \struck{Donc}
\[
\operatorname{Tr} u\, u^{\mathrm{ad}(\varphi,\psi)} > 0 \text{ pour } \forall\, u \neq 0
\]
\[
\Updownarrow
\]
\[
\operatorname{Tr} u\, u^{*} \alpha > 0 \text{ pour } \forall\, u \neq 0
\qquad (= \operatorname{Tr} u^{*} \alpha u)
\]
De plus, cette condition indique que \add{$\psi$ est de Weil, i.e.}
\[
\operatorname{Tr} u\, u^{\mathrm{ad}(\psi,\psi)} > 0 \text{ pour } u \neq 0,
\quad \text{i.e.}\quad
\operatorname{Tr} u\, \alpha^{-1} u^{*} \alpha > 0 \text{ pour } u \neq 0.
\]
\note{La ligne « $= \operatorname{Tr} u^{*}\alpha u$ » est écrite sous la
précédente, reliée par un signe $\|$ ; entre les deux conditions encadrées,
un signe qui se lit comme une double flèche verticale. Une ligne intercalée
après « i.e. » est \ill{}.}
Donc les formes de Weil sur $M$ compatibles : $\varphi$ correspondent aux
\struck{\ill{}} $\alpha \in \mathrm{End}(M)$ tels que \add{(de Weil)}
$\operatorname{Tr} u^{*} u \alpha > 0$ pour $u \neq 0$.
\note{Une insertion cerclée, rattachée ici : « Donc, dans l'ens.\ des formes
(de Weil) sur $M$, la relation de compatibilité est une relation
d'équivalence. »}
\marginal{Conséquence : si $\varphi$ et $\psi$ sont de Weil compatibles sur
$M$, \ill{} $\varphi \oplus \psi$ est de Weil, compatible avec $\varphi$ et
$\psi$}
\note{Note marginale écrite en biais au bas de la marge gauche ; deux mots
au milieu ne se lisent pas.}

6) Soient $(M_i, \varphi_i)$, \struck{des} avec
$\varphi_i \colon M_i \otimes M_i \to L$, $\varphi_i$
$\varepsilon$-symétriques de Weil ($\varepsilon = \pm 1$ fixé). Supposons que
pour $i, j$, $\varphi_i$ et $\varphi_j$ compatibles. Alors

\page{5}
cette condition est encore vérifiée si on stabilise en prenant des sommes
directes des formes et des formes \emph{induites}, sur des sous-objets.
\note{Insertion en haut de page, reliée à « sommes directes » : « de $\sum_I$,
$I$ fini, les formes sur $\bigoplus M_i$ qui sont compatibles aux $\varphi_i$
sont celles qui sont compatibles avec $\bigoplus \varphi_i$. Si $\varphi$ est
de Weil sur $M$, \ill{} ». Une note marginale cerclée, en biais : « \ill{}
$\psi$ \ill{} et $N \subset M$, \ill{} deux formes sur $N$ \ill{} sont
compatibles \ill{} $\varphi|N$ \ill{} ».}

7) \struck{Appelons} Supposons $\mathcal{M}$ graduée, avec un « objet de
Tate » $K(1)$ de degré $-2$. On appelle \emph{prépolarisation} \struck{de
$\mathcal{M}$} la donnée \struck{\ill{}}, pour $\forall\, M \in \mathrm{Ob}\,
\mathcal{M}$ de degré $i$, d'une partie \emph{non vide} \struck{\ill{}}
$P(M) \subset \mathrm{Hom}(M \otimes M, K(-i))$, formée de formes
$(-1)^{i}$-sym.\ \struck{de Weil}, satisfaisant à la condition
\begin{quote}
(P1) $\varphi \in P(M)$, $\psi \in P(N)$ ($M, N \in \mathrm{Ob}\,
\mathcal{M}^{i}$), $u \colon M \to N$ implique
$\operatorname{Tr} u^{\mathrm{ad}(\varphi,\psi)} u > 0$ si $u \neq 0$.
\end{quote}
Ceci signifie que pour tout $M \in \mathrm{Ob}\,\mathcal{M}^{i}$ ($i$ fixé)
les $\varphi \in P(M)$ sont mutuellement compatibles.
\marginal{NB Si $\exists$ prépolarisation $P$, alors $\mathcal{M}$ est
semi-simple}
Si $P$ et $P'$ sont des prépolarisations, alors
$P \subset P' \overset{\text{déf}}{\Longleftrightarrow}$ $\forall\, M$ de
degré $i$, on a $P(M) \subset P'(M)$. Notion de prépolarisation maximale.
\struck{Question : y a-t-il \ill{}}

\textbf{Proposition} Toute prépolarisation $P$ est contenue dans une unique
prépolarisation maximale $\widehat{P}$. On a $\widehat{P}(M) =$ ensemble des
formes de Weil sur $M$ compatibles aux formes de $P(M)$.

\emph{Cor.} a) Si $\varphi \in \widehat{P}(M)$, $\psi \in \widehat{P}(N)$,
alors $\varphi \oplus \psi \in \widehat{P}(M \oplus N)$ ($M$ et $N$ de même
degré $i$)

b) Si $\varphi \in \widehat{P}(M)$ \struck{\ill{}} $N \subset M$, alors
$\varphi|N \in \widehat{P}(N)$

c) Si $\varphi \in \widehat{P}(M)$, alors
$\check{\varphi} \in \widehat{P}(\check{M}(-i))$
\note{Dans c), l'argument de $\widehat{P}$ est surchargé ; la lecture
$\check{M}(-i)$ est douteuse.}

\page{6}
8) \struck{Pour} La donnée d'une prépolarisation $P$ se décompose en des
composantes $P^{i}$ ($i \in \mathbb{Z}$). Elles sont a priori tout à fait
indépendantes. Une prépolarisation est dite \emph{admissible} si elle
satisfait les conditions
\begin{quote}
(P2) $\varphi \in P(M)$, $\psi \in P(N)$, $\chi \in P(M \otimes N)
\Longrightarrow \chi$ compatible avec $\varphi \otimes \psi$, i.e.\
$\in \widehat{P}(M \otimes N)$

(P3) $\forall\, i$, l'iso can $K(i) \otimes K(i) \to K(2i)$ \struck{\ill{}} est
dans $\widehat{P}(K(i))$.
\end{quote}
\note{Dans (P3), l'exposant du membre de droite se lit mal.}
En présence de (P2), (P3) équivaut à :
\struck{(P3)$_0$ l'iso can $K(1) \otimes K(1) \to K(2)$ est dans
$\widehat{P}(K(1))$}

(P$'$3)$_1$ Si $\varphi \in P(M)$, alors $\varphi(i) \in \widehat{P}(M(i))$.
\struck{(et il \uncertain{suffit} de le supposer pour $i = 1$)}

Une \emph{polarisation} est une prépolarisation maximale et admissible.

NB Si $P$ est une prépolarisation, elle est admissible ssi $\widehat{P}$ est
une polarisation. Ceci provient au fond de ceci : si
$\varphi \in \mathrm{Hom}(M \otimes M, L)$,
$\varphi' \in \mathrm{Hom}(M' \otimes M', L')$, \uncertain{$\psi$, $\psi'$},
$\varphi \otimes \varphi' \in \mathrm{Hom}((M \otimes M') \otimes (M \otimes M'),
L \otimes L')$ ($L$, $L'$ inversibles), et si $\varphi_1$ comp.\ :
$\varphi$ et $\varphi'_1$ comp.\ : $\varphi'$, sont de Weil, alors
$\varphi_1 \otimes \varphi'_1$ compatible : $\varphi \otimes \varphi'$.
\marginal{vérifier}
\note{Le mot marginal est traversé d'un trait, peut-être une coche.}

9) Soit $P$ une polarisation. Alors on a pour les formes de $P$ :
\begin{enumerate}
\item[a)] stabilité par $\varphi \oplus \psi$
\item[a$'$)] stabilité par $\varphi + \psi$ [résulte de a) et b)]
\item[b)] stabilité par $\varphi|N$
\item[c)] stabilité par $\varphi \otimes \psi$
\item[e)] stabilité par $\check{\varphi}$
\item[d)] stabilité par $\varphi(i)$
\item[f)] \struck{Contient} \uncertain{les isos} $K(i) \otimes K(i) \simeq
K(2i)$ \struck{\ill{}} et \uncertain{leurs multiples} par $\lambda \in K$,
avec $\lambda > 0$
\end{enumerate}
\note{a$'$) est une insertion cerclée entre a) et b) ; e) est inséré, cerclé,
entre c) et d), dans cet ordre. Au bas de la page à gauche, en petits
caractères : « a) résulte \struck{\ill{}} de b) c) d) e) f) \ill{} (NB \ill{}
a) en résulte) ».}

NB Les conditions b) c) e) f) impliquent déjà que $P$ est une
prépolarisation. Les $P$ maximaux satisfaisant b) c) e) f) sont les
prépolarisations maximales satisfaisant (P2). Donc les polarisations
\uncertain{sont} \ill{}

\page{7}
10) Soit $\mathcal{M}_{\alpha}$ un système \struck{\ill{}} \uncertain{de
sous-$\otimes$-catégories de} $\mathcal{M}$, tel que avec $K(1)$ et $K(-1)$
ils engendrent $\mathcal{M}$ comme $\otimes$-catégorie. Alors une polarisation
est connue quand, pour $\forall\, \alpha$, on connaît les formes de
polarisation \uncertain{$P_{\alpha}$} des $\mathcal{M}_{\alpha}$. On peut
supposer d'ailleurs que les $\mathcal{M}_{\alpha}$ sont \uncertain{formés des
degrés $0$ et $1$}, \ill{} \uncertain{p.\ ex.\ si $\mathcal{M}$ est
$\otimes$-engendré par $\mathcal{M}^{0}$ et $\mathcal{M}^{1}$} ! \struck{\ill{}}
\note{Ce paragraphe et le suivant sont d'une écriture rapide ; les formules
passent, les phrases entre elles sont en grande partie perdues. L'exposant
« $P_{\alpha}$ » est écrit au-dessus de « de », d'une graphie proche de
$\varphi_{\alpha}$.}

11) \emph{Remarque 1} \struck{Les polarisations} Pour $M$, $L$ fixés, les
formes de Weil \struck{\ill{}} $M \otimes M \to L$ \uncertain{pour $M
\otimes_K \mathbb{R}$, sont les formes de Weil de $M \otimes_K \mathbb{R}$}
\ill{}. Les \ill{} \add{qui correspondent aux formes $\varepsilon$-sym.\
\emph{non dégénérées}, et} \ill{} \uncertain{$\mathrm{Hom}(M_{\mathbb{R}}
\otimes M_{\mathbb{R}}, L_{\mathbb{R}})$} \ill{} des prépolarisations
\uncertain{maximales} \ill{} \uncertain{correspondance biunivoque}. Donc
\ill{} $\mathcal{M}$ et $\mathcal{M}_{\mathbb{R}}$. Si $P$,
$P_{\mathbb{R}}$ \ill{} $P$ est une polarisation ssi $P_{\mathbb{R}}$
\ill{}.

\emph{Remarque 2} \ill{} $\mathcal{M}^{i}$ \ill{}
$\mathcal{M}_{\mathbb{R}}^{i}$ \ill{} polarisations \ill{} de $\mathcal{M}$
\ill{} si $\mathcal{M}$ est $\otimes$-engendré par \ill{}. \struck{c) Prouver}
\ill{}. On est ainsi ramené à étudier les polarisations, des
$\otimes$-catégories \emph{sur} $\mathbb{R}$. De plus,
\marginal{\ill{}}
\note{Une note marginale de deux lignes, en biais, à hauteur de la Remarque
2, ne se lit pas.}

\emph{Remarque 3} : Soit $\mathcal{M} = \varinjlim \mathcal{M}(\alpha)$
\add{sous-$\otimes$-catégories pleines croissantes}. Alors l'ens.\
$P(\mathcal{M})$ des polarisations de $\mathcal{M}$ est \uncertain{donné par}
$P(\mathcal{M}) = \varprojlim P(\mathcal{M}(\alpha))$ \struck{\ill{}}. Donc
si les $\mathcal{M}(\alpha)$ sont \uncertain{de t.f.}, \ill{} il

\page{8}
\uncertain{s'ensuit} que $P(\mathcal{M}) \neq \emptyset \Longleftrightarrow
\forall\, \alpha$, $P(\mathcal{M}_{\alpha}) \neq \emptyset$. On est ainsi ramené
au cas des \emph{$\otimes$-catégories semi-simples de type fini sur
$\mathbb{R}$} : étudier l'ens.\ des polarisations d'une telle $\mathcal{M}$.

12) Soit $\mathcal{M}$ avec polarisation $P$ (sur $\mathbb{R}$ si on veut).
Soit $M \in \mathrm{Ob}\,\mathcal{M}^{i}_{\mathbb{C}}$, i.e.\ un objet
$M_0$ de $\mathcal{M}^{i}$, avec action de $\mathbb{C}$ dessus (i.e.\
endomorphisme $J_M$ de carré $-1$). Soit $\overline{M}$ le $\mathbb{C}$-objet
avec l'action conjuguée ($J_{\overline{M}} = -J_M = J_M^{-1}$). Une forme
$M \otimes_{\mathbb{C}} \overline{M} \xrightarrow{\;\varphi\;} \mathbb{C}(-i)
= \mathbb{R}(-i) \otimes_{\mathbb{R}} \mathbb{C}$ équivaut : une forme
$(M \otimes_{\mathbb{C}} \overline{M})_0 \to \mathbb{R}(-i)$, i.e.\ une forme
$M_0 \otimes M_0 \xrightarrow{\;\varphi_0\;} \mathbb{R}(-i)$ satisfaisant
\[
\varphi_0(J_M x, y) = \varphi_0(x, J_{\overline{M}}\, y) = -\varphi_0(x, J_M y),
\qquad
\varphi_0(J_M x, J_M y) = \varphi_0(x, y).
\]
\note{Sous l'indice $0$ de $(M \otimes_{\mathbb{C}} \overline{M})_0$ :
« partie réelle », et une formule $1 - J \otimes \overline{J}$ ; sous la
première égalité, un « $\|$ » et un « $\mathbb{R}$ » dont le sens n'est pas
clair.}
Disons que $\varphi$ est \emph{hermitienne} resp.\ \emph{antihermitienne}
\uncertain{si} $\varphi_0$ sont sym.\ resp.\ alternée. [\ill{} hermitienne
\ill{} symétrique \ill{} correspond \ill{} $\varphi_0$]
\note{Le crochet, de deux lignes, est presque entièrement perdu.}
On dit que $\varphi$ est une polarisation (pour $P$) \struck{\ill{}} si
$\varphi_0$ \struck{\ill{}} en est une (pour $P$). Ceci implique $\varphi$
non dégénérée, i.e.\ $\tilde{\varphi} \colon \check{M} \simeq \overline{M}(-i)$

\page{9}
\ill{} une forme de polarisation \add{$P$} sur $M$, i.e.\ \struck{\ill{}} une
polarisation $\varphi_0$ de $M_0$ \emph{invariante} par $J_M$. En effet, si
$\psi_0$ est une polarisation de $M_0$, on pose
\[
\varphi_0(x,y) = \psi_0(x,y) + \psi_0(J_M x, J_M y)
\]
d'où
\[
\varphi_0(J_M x, J_M y) = \psi_0(J_M x, J_M y) + \psi_0(-x, -y) = \varphi_0(x,y).
\]
Par suite, on a : $\check{M} \simeq \overline{M}(-i)$ pour $\forall\, M \in
\mathrm{Ob}\,\mathcal{M}^{i}_{\mathbb{C}}$.)
\note{Le premier mot de la page, un $K$ ou un $Il$, ne se lit pas ; la
parenthèse fermante finale est sur la page.}

13) On voit donc que les \uncertain{implications} :
\begin{quote}
$\alpha$) $\mathcal{M}$ admet une polarisation \add{sur $\mathbb{R}$}

$\Downarrow$

$\beta$) $\forall\, M \in \mathrm{Ob}\,\mathcal{M}^{i}_{\mathbb{C}}$,
$\check{M} \simeq \overline{M}(-i)$, $\mathcal{M}$ ½-simple

$\Downarrow$

$\gamma$) $\forall\, M \in \mathrm{Ob}\,\mathcal{M}^{i}_{\mathbb{R}}$,
$\check{M} \simeq M(-i)$, $\mathcal{M}$ ½-simple
\end{quote}
\marginal{Lemme}
($\beta \Longrightarrow \gamma$ provient du fait que $M_{\mathbb{C}} \simeq
N_{\mathbb{C}}$ implique $M \simeq N$).
\note{Les mentions « $\mathcal{M}$ ½-simple » sont cerclées ; un « ? » dans
la marge gauche à hauteur de la première implication.}
Les implications précédentes \uncertain{ne se renversent pas} : $\mathcal{M}
= \mathrm{Rep}(G)$, $G$ groupe algébrique \struck{\ill{}} sur $\mathbb{R}$. On
peut avoir $\beta$) \struck{sans avoir $\alpha$)} [\struck{$SO(\varphi)$}
$SO(\varphi)/\text{centre}$, $\varphi$ forme \uncertain{en} un nombre pair de
variables, \uncertain{avec} un nombre impair de carrés négatifs], et
$\gamma$) sans avoir $\beta$) [\ill{}].
\note{Le second crochet, d'une ligne, est perdu.}

\page{10}
14) Supposons $\mathcal{M}$ semi-simple \add{sur $K \subset \mathbb{R}$},
engendrée \struck{\ill{}} \add{et $K(1)$, $K(-1)$} par des $M_1, \ldots, M_r$
de degrés $d_1, \ldots, d_r$, \struck{(Pour tt $i$,} \struck{\ill{}} et de
formes $\varepsilon^{d_i}$-sym.\ sur $M_i$, qui soit \struck{\ill{}} formes de
Weil. \struck{\ill{}} Pour qu'il existe une polarisation de $\mathcal{M}$
\struck{définie} pour laquelle les $\varphi_i$ soient des formes de
polarisation, il faut qu'on ait les conditions suivantes :
\note{Le sujet de « qui soit » manque sur la page ; les $\varphi_i$ ne sont
nommées qu'à la phrase suivante.}

a) Pour deux systèmes d'entiers $m = (m_1, \ldots, m_r)$,
$m' = (m'_1, \ldots, m'_r) \geqslant 0$, avec
$m_{r+1} + \sum_{1}^{r} m_i d_i = m'_{r+1} + \sum m'_i d_i$, posant
\[
M_m = \Bigl(\bigotimes_{1 \leqslant i \leqslant r} \bigl(\otimes^{m_i} M_i\bigr)\Bigr)(m_{r+1}),
\qquad
\varphi_m = \Bigl(\bigotimes_{1 \leqslant i \leqslant r} \bigl(\otimes^{m_i} \varphi_i\bigr)\Bigr)(m_{r+1}),
\]
on a, pour tout $u \in \mathrm{Hom}(M_m, M_{m'})$, la relation
\[
\operatorname{Tr} u^{*(\varphi_m, \varphi_{m'})} u \geqslant 0.
\]
\note{Au-dessus des deux sommes, l'annotation « $\delta(m)$ » ; le signe de la
relation se lit $\geqslant$, non $>$.}

Je dis que cette condition est aussi suffisante. \struck{Supposons réalisée
en effet} \uncertain{Elle étant remplie}, si $M \in \mathrm{Ob}\,
\mathcal{M}^{n}$, $M \neq 0$, $\exists\, m$ avec un $\alpha \colon M
\hookrightarrow M_m$ ($n = \struck{\ill{}} \sum m_i d_i$), soit $P(M)$ l'ens.\
des formes qui sont $\sim \varphi_m|M$. \uncertain{Prouvons qu'il est
indépendant de $m$, $\alpha$.} En effet, si \ill{} $M_{m'}$, comme
$\varphi_m$ et $\varphi_{m'}$ sont \ill{} et $\alpha' \colon M
\hookrightarrow M_{m'}$, il en est de même, compatibles par hyp.,
\struck{\ill{}} $\varphi_m|M$ et $\varphi_{m'}|M$. On voit que ces $P(M)$
forment une prépolarisation \struck{maximale}, \ill{} : \uncertain{Cependant à
\emph{chacune} des $\varphi_i$}. Si $\varphi \in P(M)$, $\varphi' \in P(M')$,
$\varphi \sim \varphi_m|M$, $\varphi' \sim \varphi_{m'}|M'$, donc
$\varphi \otimes \varphi' \sim \varphi_m \otimes \varphi_{m'}|M \otimes M'
= \varphi_{(m+m')}|M \otimes M'$,
\note{« $\sim$ » est ici la compatibilité. Dans la phrase « Prouvons… », le
mot lu « indépendant » est écrit vite ; la ligne suivante est très
abrégée.}

\page{11}
donc $\varphi \otimes \varphi' \in P(M \otimes M')$. Enfin, supposons
$\varphi \in P(M)$, i.e.\ $\varphi \sim \varphi_m|M$, alors
$\varphi(i) \sim \varphi_m(i)|M(i) = \varphi_{(m_1, \ldots, m_r, m_{r+1}-i)}|M(i)$,
donc $\varphi(i) \in P(M(i))$.

15) Soit $G$ groupe algébrique sur $\mathbb{R}$, \struck{et} avec un
homomorphisme
\[
G_m \xrightarrow{\;i\;} G \xrightarrow{\;\varepsilon\;} G_m,
\]
$\eta = i(-1)$, et soit $C \in G(\mathbb{R})$ avec $C^{2} = \eta$. On prend
$\mathcal{M} = \mathrm{Rep}(G)$. \struck{1) Soit $M$}

a) Si $M \in \mathrm{Ob}\,\mathcal{M}^{i}$, i.e.\ (rep.\ de \struck{degré}
poids $i$, \ill{}), et si $\varphi \colon M \otimes M \to \mathbb{R}(-i)$ est
telle que $\varphi(x, Cy)$ est symétrique et \struck{une} forme définie
positive, on dit que $\varphi$ est de Weil, et si $\psi \colon N \otimes N
\to \mathbb{R}(-i)$, $\psi(x, Cy)$ déf.\ positive, alors $\varphi$ et $\psi$
sont \emph{compatibles}.
\marginal{$\varphi$ est dite une forme de $C$-polarisation}
\marginal{\uncertain{vérifions d'abord} que $\varphi$ est \struck{\ill{}}
$(-1)^{i}$-sym.}
\note{Une insertion au-dessus de « poids $i$ », de deux mots, ne se lit
pas.}
En effet, soit $u \colon M \to N$ compatible avec l'action de $G$.
Considérons $u^{*(\varphi,\psi)}$, et comparons avec
$u^{*(\varphi^{C},\psi^{C})}$ ($\varphi^{C}(x,y) = \varphi(x, Cy)$). On a
\[
u^{*(\varphi^{C},\psi^{C})} = C_M^{-1}\, u^{*(\varphi,\psi)}\, C_N,
\quad\text{donc}\quad
u^{*(\varphi,\psi)} = \underbrace{C_M\, u^{*(\varphi^{C},\psi^{C})}\, C_N^{-1}}_{u'},
\]
donc
\[
\operatorname{Tr} u\, u^{*(\varphi,\psi)}
= \operatorname{Tr} u\, C_M\, u'\, C_N^{-1}
= \operatorname{Tr} C_N^{-1} u\, C_M\, u'
= \operatorname{Tr} u\, u' \geqslant 0.
\]
\note{L'accolade « $u'$ » est sous $C_M u^{*(\varphi^C,\psi^C)} C_N^{-1}$ ;
une seconde accolade sous $C_N^{-1} u C_M$ porte une annotation lue « $u$ ».
Dans la marge gauche, en biais : « D'après la condition de positivité
\uncertain{(Cor.\ \ill{})} : $\operatorname{Tr} u\, u^{*(\varphi,\psi)} =
\operatorname{Tr} u\, u^{*(\varphi^{C},\psi^{C})}$ ».}
\struck{\ill{}}
\note{Suivent trois lignes biffées, dont on lit « orthogonalement pour les
formes », « $G(\mathbb{R})$ » et « orthogonal », puis à gauche un bloc
biffé de quatre lignes où l'on lit « $C_M$ et $C_N$ », « $\varphi^{C}$,
$\psi^{C}$ », « $\mathrm{Cent}(C)$ » et « $C_N^{-1} = C$ ».}

b) Si $\varphi$, $\psi$ sont des $C$-polarisations, $\varphi \otimes \psi$
aussi, \add{pourvu que $\varepsilon(C) = 1$ (et non $-1$ !)} ; si $\varphi$
est une $C$-polarisation, alors $\varphi(i)$ aussi ; la forme standard
$\mathbb{R} \otimes \mathbb{R} \to \mathbb{R}$ est une $C$-polarisation.

\page{12}
\ill{} que une $C$-polarisation de \ill{}, les formes $\varphi'$ sur \ill{}
$= M$ qui sont les $C$-polarisations, sont \emph{exactement} celles qui sont
compatibles à $\varphi$.

c) Donc : Supposons $\varepsilon(C) = 1$. Soit, pour tt $M$ de \add{poids
$i$}, $P(M)$ l'ens.\ des $C$-polarisations. Alors $M \mapsto P(M)$ est une
polarisation de $\mathcal{M}$, et \uncertain{ceci prouve} que les $P(M)$
\uncertain{peuvent en fait être} tous $\neq \emptyset$.

16) \emph{Questions} : Toute polarisation en \struck{\ill{}} $\mathcal{M} =
\mathrm{Rep}(G)$ provient-elle \add{se dérive} par un $C$ comme dessus
($C^{2} = \eta$, $\varepsilon(C) = 1$) ? Est-ce vrai si $G$ connexe ? ou $G'$
\add{$= \operatorname{Ker} \varepsilon$} \uncertain{compact} ?
\note{« $= \operatorname{Ker}\varepsilon$ » est cerclé au-dessus de $G'$ ;
à côté, cerclé : « Cas $G'$ compact ? ».}

NB Un tel $C$ \add{est unique mod centre $Z$} \struck{\ill{}} est
\uncertain{déterminé} mod le groupe des \struck{\ill{}} \ill{}
\uncertain{$G$ connexe ?} conjugaison, \struck{\ill{}} \ill{} $Z$. Donc
\struck{bien} l'indétermination de $C$ dans ce \add{cas} est modulo points
d'ordre $2$ de $Z(\mathbb{R})$ et conjugaison.

Re NB Pour qu'il existe un $C$ comme dessus fournissant une polarisation, il
\uncertain{faut et suffit} \add{\ill{}} que $G' = \operatorname{Ker}
\varepsilon$ muni de $\eta = i(-1)$, \uncertain{soit tel} qu'il existe un
$C \in G'(\mathbb{R})$ tel que \add{\ill{}} rep.\ \add{de $G'$ sur lesquelles
$\eta$ \uncertain{opère par} $\pm 1$} \struck{\ill{}} une $C$-polarisation,
i.e.\ une forme \add{inv.} $\varphi$ telle que $\varphi(x, Cy)$ soit sym.\
déf $> 0$. On dit que $G'$ est \emph{hodgien} s'il existe un tel $C$ (avec
$C^{4} = 1$) \struck{d'ordre $4$} ($\eta = C^{2}$ n'est pas fixé a priori).
\note{Le Re NB est chargé d'insertions interlinéaires, dont deux ne se
lisent pas ; « faut et suffit » est écrit « f.\ et s. ».}

\page{13}
17) Soit $C \in G'(\mathbb{R})$, avec $C^{2} = \eta$ ($= i(-1)$),
définissant une polarisation de \add{$K$} $\mathcal{M} = \mathrm{Rep}(G)$, ---
donc $G$ hodgien ! Alors $\mathrm{Centr}_{G'}(C)(\mathbb{R})$ est un compact
maximal de $G'(\mathbb{R})$. En effet, on \struck{\ill{}} \struck{$G' \subset
\mathrm{Aut}_{\mathbb{R}}(M) \times \mathrm{Aut}_{\mathbb{R}}(M')$}
\struck{\ill{}} \add{a} $G' \hookrightarrow \prod_{\alpha}
\mathrm{Aut}_{\mathbb{R}}(M_{\alpha})$, pour des rep.\ convenables
\add{(en nb de $2$ si on veut)} de $G'$, et \add{avec} $\eta$
\uncertain{opérant} par $\pm 1$. Si $\varphi_{\alpha}$ est une forme de
$C$-polarisation sur $M_{\alpha}$, alors \uncertain{le gr.}
$\mathrm{Centr}_{G'}(C)$ \uncertain{laisse invariant} $\varphi^{C}_{\alpha}$,
donc \struck{$\mathrm{Cent}$} $K \subset \underbrace{\prod_{\alpha}
\mathrm{Aut}_{\mathbb{R}}(M_{\alpha}, \varphi^{C}_{\alpha})}_{\text{groupes
orthogonaux compacts}}$, donc $K$ est compact \ill{}. Prouvons
\add{\uncertain{compacte}} qu'il est \struck{\ill{}} \uncertain{maximal}.
\struck{\ill{}} Soit \struck{donc} $K' \subset G'(\mathbb{R})$, $K' \supset
K(\mathbb{R})$, \add{\ill{}} $K'$ compact maximal. \struck{\ill{}} On sait
que $K'$ est algébrique, $K' = \mathbf{K}'(\mathbb{R})$, $\mathbf{K}'$
\uncertain{sous-groupe} algébrique de $G'$ \add{\ill{}} \struck{\ill{}} $G'$
contient $C$. Quitte à remplacer $G'$ par $K'$, on peut supposer $G'$
compact, et il faut prouver qu'alors $C$ est central. Prenons le d'abord
dans $G'/\{1, \eta\}$, i.e.\ supposons donc \uncertain{$\eta = 1$}. Soit une
rep.\ fidèle de $G'$ dans un $V$, avec forme de $C$-pol.\ $\varphi$ sur $V$,
\struck{Alors \ill{}} \struck{$C|V^{+} = \mathrm{id}$, $C|V^{-} = -\mathrm{id}$}
\add{alors} $V = V^{+} + V^{-}$ et $(\varphi|V^{+}) > 0$,
$(\varphi|V^{-}) < 0$. \struck{On a} \add{Comme} $G'$ est un sous-groupe de
$O(\varphi)$, on \uncertain{peut} \add{par conjugaison \uncertain{par un
élément de $O(\varphi)(\mathbb{R})$}, l'envoyer dans le compact maximal de
$O(\varphi)$} $O(V^{+}) \times O(V^{-})$, donc on peut supposer $G \subset
O(V^{+}) \times O(V^{-})$. Donc $C \in G(\mathbb{R})$ est de la forme
$C^{+} \oplus C^{-}$, \struck{\ill{}} avec $\varphi^{+}(x, C^{+} y)$ et
\note{« $K$ » n'est pas défini sur la page ; c'est le centralisateur qui
vient d'être nommé. Sous « $(\varphi|V^{-}) < 0$ » : « $= -\varphi^{-}$ ».
Le passage sur le compact maximal de $O(\varphi)$ est écrit dans la marge
gauche, en petits caractères, sous une double barre.}

\page{14}
$\varphi^{-}(x, C^{-} y)$ sur $V^{+}$ resp.\ $V^{-}$ \add{\ill{}} définies
positives, \struck{\ill{}} $C^{+}$ et $C^{-}$ \add{sont des} involutions
orthogonales, \ill{} on \uncertain{conclut} que $C^{+} = \mathrm{id}$,
$C^{-} = -\mathrm{id}$, donc $C$ \add{est} dans \struck{$O(\varphi)$}
\add{le centre de} $O(V^{+}) \times O(V^{-})$ donc, dans le centre de $G$.

Cas $\eta \neq 1$. Alors \struck{\ill{}}
\note{Trois lignes biffées, dont on lit « un morphisme $G$ », « $\{1, \eta\}$ »,
« $gCg^{-1} \in \{1, \eta\}$ », « nulle ».}
Alors $gCg^{-1} \in \{C, \eta C\}$, donc $G'$ \struck{\ill{}} \struck{pour
$g \in G(\mathbb{R})$}.
\note{Un long passage encadré d'une accolade dans la marge gauche, de sept
lignes, biffé en partie par un trait oblique, où l'on lit : « Soit une
représentation \ill{} il suffit de prouver que \ill{} central. Donc \ill{}
$\eta$ opérant par $-\mathrm{id}$. Donc \ill{} considérons \ill{} $G \subset
\mathrm{Aut}_{\mathbb{R}}(M)$, \ill{} l'image, $C$ devient \ill{} Donc $C$
\uncertain{définit} sur $M$ \ill{} \struck{$G \subset
\mathrm{Aut}_{\mathbb{C}}(M)$} Les formes \ill{} satisfont à \ill{}
$\varphi(\cdot, \cdot\, C)$ ».}
\uncertain{Considérons} \ill{} soit $G''$ le stabilisateur \add{de $C$},
d'indice $1$ ou $2$, \struck{$G''$ et \ill{}} ($\varepsilon \colon G \to
\mu_2$)
\[
g C g^{-1} = \varepsilon(g)\, C,
\qquad
\varphi^{C}(gx, gy) = \varphi(gx, Cgy) = \varepsilon(g)\, \varphi(gx, gCy)
= \varepsilon(g)\, \varphi^{C}(x, y).
\]
\struck{Donc $G' \subset Sp(\varphi^{G})$} Si on avait $\varepsilon(g) = -1$,
on aurait $\varphi^{C}(gx, gx) = -\varphi^{C}(x, x)$, et si $x \neq 0$ les
deux membres seraient resp.\ $> 0$ et $< 0$, absurde.
\note{Le $\varepsilon$ de cette page, caractère de $G$ à valeurs dans
$\mu_2$, n'est pas celui de 15) ; la page réemploie la lettre.}

\page{15}
18) Supposons que $C$, $C'$ \add{($C^{2} = C'^{2} = \eta$)} définissent des
polarisations de $\mathcal{M}$ \add{$= \mathrm{Rep}(G)$}. Quand celles-ci
sont-elles égales ? Il \emph{suffit} qu'on ait $C' = gCg^{-1}$, avec $g \in
G(\mathbb{R})$, car alors pour $\varphi$ forme invariante \ill{}, on aura
\[
\varphi^{C'}(x, y) = \varphi(x, gCg^{-1} y) = \varphi(g^{-1}x, Cg^{-1}y)
= \varphi^{C}(g^{-1}x, g^{-1}y).
\]
Cette condition est-elle aussi nécessaire ? \uncertain{Question} \ill{}
\add{$G'$ connexe} \add{$G'(\mathbb{R})^{\circ}$ !} Tout d'abord,
\struck{\ill{}} \ill{} par un élément de $G'(\mathbb{R})^{\circ}$ \ill{}
\uncertain{des groupes compacts maximaux} \ill{}, il existe $g \in
G'(\mathbb{R})^{\circ}$ tel que
\[
g C g^{-1} C'^{-1} \in Z'(\mathbb{R})
\]
($Z' =$ centre de $G'$) [NB $(G'/Z')(\mathbb{R})^{\circ} \longleftarrow
G'(\mathbb{R})^{\circ}$ est surjectif !]. Donc on aura
\[
C' = z\, g C g^{-1}, \qquad g \in G'(\mathbb{R})^{\circ},\ z \in Z'(\mathbb{R}).
\]
D'ailleurs $z^{2} = 1$ i.e.\ $z \in {}_2Z'(\mathbb{R})$. Quitte à remplacer
$C'$ par \ill{}, on aura $C' = zC$, $z \in {}_2Z(\mathbb{R})$, \uncertain{et}
on va prouver que si $z \neq 1$, $C$ et $C'$ donnent des polarisations
diff. Or si $z \neq 1$, il y a une rep.\ \uncertain{fidèle} $V \neq 0$ de
$G'$ où $z$ opère par $-\mathrm{id}_V$. Si $\varphi$ est une forme de
$C$-polarisation, $\varphi^{C} > 0$, alors
\[
\varphi^{C'} = \varphi^{zC} = -\varphi^{C} < 0,
\]
absurde ! Donc on trouve que \struck{\ill{}} \emph{l'ensemble des
polarisations} sur $\mathcal{M} = \mathrm{Rep}(G)$ qui sont définissables par
un $C$ (i.e.\ qui sont « hodgiennes ») forme un \emph{torseur} sous
${}_2Z(\mathbb{R})$ : Si $P$ est une polarisation \add{hodgienne} et $z \in
{}_2Z(\mathbb{R})$, soit $zP$ \uncertain{est} la polarisation
\note{Dans la marge gauche, cerclé, à hauteur de la formule pour
$\varphi^{C'}$ : « \struck{\ill{}} NB $C'$ commute \uncertain{avec} $C$
\uncertain{et} \ill{} l'\uncertain{existence} \ill{} $C$ \ill{} $C' = \eta$ ».
Le passage « Tout d'abord … compacts maximaux » est de trois lignes très
abrégées.}

\page{16}
telle que les formes $\varphi$ sur $M$ qui sont des $(zP)$-polarisations,
sont celles pour lesquelles $\varphi^{z}$ (défini par $\varphi^{z}(x,y) =
\varphi(x, zy)$) sont \struck{définies} des $P$-polarisations.

19) Le \emph{cas général}, $\mathcal{M}$ sur $K \subset \mathbb{R}$,
\add{$\otimes$-cat.\ quelconque}, si $\varphi$ est une forme $M \otimes M
\to L$ \add{\ill{}}, \struck{\ill{}} et $z \in {}_2\mathrm{Aut}_{\otimes}
(\mathrm{id}_{\mathcal{M}}) =: {}_2Z(K)$ ($Z =$ centre des \uncertain{groupes
de \ill{} de $\mathcal{M}$}), on définit $\varphi^{z}$ sur $M \times M$ par
$\varphi^{z}(x,y) = \varphi(x, zy)$. On \uncertain{trouve} alors, si
$\varphi^{z}$ et $\psi$ (sur $N$) sont \uncertain{$\varepsilon$-sym.}
($\varepsilon = \pm 1$), et si $u \colon M \to N$, $u^{*(\varphi,\psi)}
\colon N \to M$, alors
\[
u^{*(\varphi^{z}, \psi^{z})} = z_M^{-1}\, u^{*(\varphi,\psi)}\, z_N
= u^{*(\varphi,\psi)}.
\]
Donc : \struck{$\varphi$ et $\psi$ sont} $\varphi$ de Weil
$\Longleftrightarrow$ $\varphi^{z}$ de Weil, $\varphi$ et $\psi$ de Weil
comp.\ $\Longleftrightarrow$ idem pour $\varphi^{z}$ et $\psi^{z}$.
\struck{D'ailleurs on constate} \struck{\ill{}}
$\varphi^{z} \otimes \psi^{z} = (\varphi \otimes \psi)^{z}$. Donc si $z$
opère par $\mathrm{id}$ sur $K(1)$ (i.e.\ $z \in {}_2Z'(K)$)
\struck{pour que} on voit que si $P$ est une polarisation, alors
\[
{}^{z}(P(M)) = \{ {}^{z}\varphi \mid \varphi \in P(M) \}
\qquad (M \in \mathrm{Ob}\,\mathcal{M}^{i},\ i \in \mathbb{Z})
\]
forme une polarisation ${}^{z}P$ sur $\mathcal{M}$. On voit tout de suite
que cette dernière \uncertain{ne peut} \ill{} être égale à $P$ que si
$z = 1$. Ainsi, ${}_2Z'(\mathbb{R})$ opère librement sur
$\mathrm{Pol}(\mathcal{M})$. \emph{Questions} : $\mathrm{Pol}(\mathcal{M})$
est-il un pseudo-torseur sous ${}_2Z(\mathbb{R})$ ?
\marginal{la condition \ill{} $z^{2} = 1$ sert \ill{} que $\varphi$
$\varepsilon$-sym.\ \uncertain{$\Longrightarrow$} $\varphi^{z}$
$\varepsilon$-sym.}
\note{La note marginale est cerclée, en biais, à hauteur de la formule pour
$u^{*(\varphi^z,\psi^z)}$. Le passage de $Z'(K)$ à $Z'(\mathbb{R})$ puis
$Z(\mathbb{R})$ dans les dernières lignes est sur la page. Un « ? » dans la
marge gauche, avec une double barre, à hauteur de la question finale. La
notation passe de $\varphi^{z}$ à ${}^{z}\varphi$ au milieu de la page.}

\section*{Notes anciennes (pages 18 et 20)}

\note{Le titre est de sa main, en haut à droite de la page 18, dans un
cadre : « Notes anciennes ». Ces deux feuillets, d'un papier et d'une encre
différents, portent une notation plus ancienne : le motif $M$ de poids $p$,
$\mathbb{Q}(p)$, $M^{*} = \check{M}(p)$, $\mathcal{Q}(M) = B(M, M ;
\mathbb{Q}(p))$ l'espace des formes bilinéaires. Le glyphe dans
« $\mathbb{Q}(p)$ » et « $\check{M}(p)$ » se lit entre $p$ et $\varphi$ ; la
lecture $p$ suit « poids $p$ » de la première ligne.}

\page{18}
$\varphi$ forme symétrique resp.\ \uncertain{alternée} sur le motif $M$ de
poids $p$, $\varphi$ à valeurs dans $\mathbb{Q}(p)$, $\varphi$ non dégénérée

$\varphi$ définit un isomorphisme
\[
\mathrm{End}(M) \simeq \struck{\ill{}} \mathcal{Q}(M)
\]
en associant à toute forme $\psi$ l'endomorphisme
$\tilde{\varphi}^{-1}\tilde{\psi}$ de $M$, et à tout endomorphisme $u$ de $M$
la forme $\varphi \circ (u \otimes \mathrm{id})$, \struck{donc}
\uncertain{correspondant à} $\tilde{\varphi} \circ u$.
\[
M \underset{\tilde{\psi}}{\overset{\tilde{\varphi}}{\rightrightarrows}} M^{*} = \check{M}(p),
\qquad
\varphi(x,y) = \struck{\ill{}} (y, \tilde{\varphi}(x))
\]
\note{Deux termes biffés, l'un après « $\simeq$ », l'autre avant « $(y,
\tilde{\varphi}(x))$ », sont illisibles sous la rature.}
Alors le \emph{symétrique} \struck{\ill{}} de $\mathcal{Q}(M)$ devient une
application \uncertain{de} $\mathrm{End}(M)$ en lui-même, notée
$u \mapsto u^{\varphi}$, qui n'est autre que l'\emph{adjoint} relativement à
$\varphi$,

\page{20}
\[
\mathcal{Q}(M) = B(M, M ; \mathbb{Q}(p)) \simeq \mathrm{Hom}(M, M^{*})
\qquad \text{où } M^{*} = \check{M}(p)
\]
par $\langle \tilde{\varphi}(x), y\rangle = \varphi(x,y)$
\note{Sous $\tilde{\varphi}(x)$ : « $\in M^{*}$ » ; sous $y$ : « $\in M$ ».}

Si $\tilde{\varphi}$ est un isom., on trouve
\[
\mathrm{End}(M) \xrightarrow[\;\sim\;]{\;\Theta_{\varphi}\;} \struck{\ill{}}\ \mathcal{Q}(M) \simeq \mathrm{Hom}(M, M^{*})
\]
en associant à $u$
\[
\Theta_{\varphi}(u) = \varphi \circ (u \otimes 1)
\qquad\text{donc}\qquad
\widetilde{\Theta_{\varphi}(u)} = \tilde{\varphi} \circ u
\]
et en sens inverse
\[
\Theta_{\varphi}^{-1}(\psi) = \tilde{\varphi}^{-1} \tilde{\psi}
\qquad
M \underset{\tilde{\psi}}{\overset{\tilde{\varphi}}{\rightrightarrows}} M^{*}
\]
La symétrie (dans $\mathcal{Q}(M)$) $\psi \mapsto \psi' = \ill{}\ {}^{s}\psi$,
devient un endomorphisme $s$ de $\mathrm{End}(M)$, \struck{\ill{}} savoir
$u \mapsto u^{\varphi}$,
\[
u^{\varphi} = \struck{\ill{}}\ \tilde{\varphi}^{-1}\bigl[{}^{s}[\varphi \circ (u \otimes 1)]\bigr]
= \tilde{\varphi}^{-1}\bigl(\varphi^{s} \circ (1 \otimes u)\bigr)
\]
de sorte que
\[
\struck{\varphi(ux, y) = \varphi(x, u^{\varphi} y)}
\]
\[
\varphi(u^{\varphi} x, y) = (-1)^{\delta} \varphi(u\struck{x}, \struck{x})
\]
\[
(-1)^{\delta}\, \Theta_{\varphi}(u)(x,y) = \Theta_{\varphi}(u)'(y,x)
\]
donc si $\varphi$ est symétrique (\add{et $p$ pair}, ou \ill{})
\[
\struck{(a)}\quad \varphi(ux, y) = \varphi(x, u^{\varphi} y)
\]
si $\varphi$ est \struck{\ill{}} antisymétrique \struck{et $p$ impair, c'est}
\[
\struck{(a)}\quad \varphi(ux, y) = \struck{\ill{}}\ \varphi(x, u^{\varphi} y)
\]
\note{Dans la définition de $u^{\varphi}$, un facteur biffé et surchargé
précède $\tilde{\varphi}^{-1}$, au-dessus duquel est écrit « $(-1)^{\delta}$
$\tilde{\varphi}^{-1}$ » ; l'exposant $\delta$ n'est pas défini sur ces deux
pages. Dans la formule pour $\varphi(u^{\varphi}x, y)$, les deux arguments
du membre de droite sont biffés d'une croix et la formule est laissée
telle quelle. Les deux dernières lignes, symétrique et antisymétrique,
aboutissent à la même égalité ; le signe biffé dans la seconde ne se lit
pas.}

\end{document}
